\chapter{Proprietà FLTL e verifica}

Le proprietà riportate di seguito sono quelle effettivamente definite
nel modello FSP. La gestione dell'allarme è incorporata nelle proprietà:
la coda sopravvive all'evacuazione, mentre una corsa interrotta può
terminare con l'evento \texttt{evacuate}.
\begin{description}
  \item[\texttt{MUTUAL\_EXCLUSION}] dopo un ingresso, prima di un
    nuovo ingresso devono avvenire \texttt{exit} oppure
    \texttt{evacuate}:
    \[
      \Box\bigl(\mathit{enter}\Rightarrow
      \neg\mathit{enter}\ \mathsf{U}\
      (\mathit{exit}\vee\mathit{evacuate})\bigr).
    \]
  \item[\texttt{EVENTUAL\_ARRIVAL}] ogni chiamata è seguita dal
    raggiungimento del piano dell'utente con le porte aperte:
    \[
      \Box\bigl(\mathit{call}\Rightarrow
      \mathsf{F}(\mathit{LIFT\_AT}(h)\land
      \mathit{DOORS\_OPEN})\bigr).
    \]
  \item[\texttt{EVENTUAL\_DELIVERY}] ogni selezione è seguita
    dall'uscita oppure dall'evacuazione:
    \[
      \Box\bigl(\mathit{select}\Rightarrow
      \mathsf{F}(\mathit{exit}\vee\mathit{evacuate})\bigr).
    \]
  \item[\texttt{STARVATION\_FREEDOM}] un utente che ha effettuato una
    chiamata e si trova nello stato di attesa raggiunge prima o poi
    \texttt{enter}:
    \[
      \Box\bigl(\mathit{WAITING}(h,u)\Rightarrow
      \mathsf{F}\mathit{enter}\bigr).
    \]
\end{description}

In LTSA la proprietà di mutua esclusione può essere espressa come
processo parametrico:
\begin{verbatim}
property MUTUAL_EXCLUSION =
    (enter[h:FLOORS][u:USER_IDX] -> IN_CABIN),
IN_CABIN =
    (exit[h:FLOORS][u:USER_IDX][t:FLOORS] ->
        MUTUAL_EXCLUSION
    |evacuate[h:FLOORS][u:USER_IDX] ->
        MUTUAL_EXCLUSION).
\end{verbatim}
Nel modello, le proprietà sono parametrizzate sugli stessi indici
\texttt{FLOORS}, \texttt{USER\_IDX} e, per la consegna,
\texttt{FLOORS} della destinazione.

\section{Esito dell'analisi}
La struttura del modello garantisce la mutua esclusione: \texttt{LIFT}
abilita \texttt{enter} soltanto nello stato \texttt{EMPTY} e porta la
cabina nello stato \texttt{FULL}. Le guardie del processo di coda
garantiscono inoltre che solo il ticket con \texttt{k == 0} possa
richiedere il servizio.

Le proprietà di vivacità sono verificate sotto l'ipotesi di scelta
fair di LTSA. \texttt{EVENTUAL\_ARRIVAL} e
\texttt{STARVATION\_FREEDOM} possono attraversare un allarme, perché
la coda viene conservata e il servizio riprende dopo
\texttt{alarm\_solved}. \texttt{EVENTUAL\_DELIVERY} ammette invece
\texttt{evacuate} come alternativa all'uscita alla destinazione.
L'output numerico della verifica LTSA non è tra gli artefatti presenti
e va prodotto eseguendo le proprietà sul composto \texttt{SYS}.